\chapter{From a law to an executable claim}\label{ch:software-claim}
A mathematical equation can fit on one line while its faithful implementation
requires many boundaries. The equation does not specify how state is frozen,
how an action is identified, which module may read another cell, or what
happens if a record is written twice. These are not decorative engineering
details. They determine whether the program computes the object named in
the theorem.

\section{Four kinds of completion}
A derived identity is mathematically complete when its statement and proof
have the required assumptions. A component is implemented when its actual
code performs its declared function. A test is completed when it has
executed and produced inspectable output. A scientific study is completed
only under its registered execution and finalization contract. A file named
\texttt{runner.py} establishes none of the last three by its name alone.

The current Gaussian direction is prospective \cite{programme}. The
repository contains historical executable engines, preserved candidate
harness code and framework objects on source-locked lineages. Their presence
does not establish an accepted Gaussian runtime. This chapter uses
``would'' and ``must'' for prospective responsibilities, not as euphemisms
for completed implementation.

\section{A module's scientific contract}
For each component, state its inputs, permitted reads, outputs, permitted
writes, refusal conditions and version identity. A value evaluator should
not choose the action. A physical checker should not improve an action's
score. A selector should not alter quantities after valuation. A logger
should not repair a failed invariant by changing history.

For example, an evaluator might receive a frozen state projection, reference
and scale declarations, and a finite accepted group. It would return the
exact field value and declared attribution record without changing physical
state. An execution adapter would later apply the exact selected group,
under separate permission. Keeping those calls separate makes it possible
to test arithmetic without accidentally running a model.

\section{Why a generic framework is still useful}
The source-inspected F2 decision \cite{framework} identifies reusable
identity, canonical serialization, typed state/action/support records,
information contracts, provenance and ledger structures. It does not say
the whole framework already fits the Gaussian programme unchanged.

In particular, its potential declaration is metadata, not a numerical
Gaussian evaluator. Its named scientific advancement routes are deliberately
guarded. Its physical service-capacity records are not earned EBU balances.
Reusing those names while changing their meanings would conceal a new
scientific model inside an apparently routine software integration.

\section{No framework replacement by default}
The smallest honest path is to reuse inspected generic pieces and add
separately versioned adapters where the active scientific meaning differs.
Wholesale copying can import historical assumptions; wholesale replacement
can discard proven provenance discipline. The planning decision is neither
``everything works already'' nor ``nothing can be reused.''

The source locks matter. Some relevant framework files live on another
lineage and are absent from the active checkout. An ignored bytecode
directory is not an authoritative source package. A later implementation
task must inspect the exact committed files it intends to reuse, rather
than infer availability from import names or an old chat summary.

\section*{Worked review question}
A function named \texttt{advance\_epoch} exists but always refuses because
its scientific execution permit is unavailable. Is epoch advancement
implemented? The refusal guard is implemented. Successful scientific
advancement is not established. Both facts belong in the status report;
neither ``there is no software'' nor ``the engine is ready'' is accurate.

\section*{Chapter recap}
The bridge from mathematics to software is a set of typed responsibilities
and verifiable effects. The bridge from software to evidence requires an
additional registered study, not a stronger adjective in the implementation
description.
